\documentclass[answers]{exam}

% Version compilable
%\usepackage{../../../mypackages}
%\usepackage{../../../macros}

% Version pour execution du script
\usepackage{../../../mypackages}
\usepackage{../../../macros}

% Pour le tableau
\usepackage{array}
\usepackage{longtable}
\usepackage{multirow}
\usepackage[table]{xcolor}

\title{Devoir en classe - Geogebra + Tracés géométriques}
\author{Nicolas Bancel}
\date{Mardi 15 Décembre 2025}

% Mise en forme des solutions en bleu
\SolutionEmphasis{\color{blue}}
\renewcommand{\solutiontitle}{\noindent}

\definecolor{lavande}{RGB}{180,190,255}

% Colonne p{} alignée à gauche et qui wrappe
\newcolumntype{L}[1]{>{\raggedright\arraybackslash}p{#1}}

\begin{document}

\textbf{Collège Lycée Suger}
\hfill
\textbf{Mathématiques} \\

\textbf{Année 2025-2026}
\hfill
\textbf{Classe de 6ème} \par

{\let\newpage\relax\maketitle}

\begin{tcolorbox}[
  %colback=blue!5!white,
  colframe=blue!90!black,
  %title=Commentaire plus détaillé,
  fonttitle=\bfseries,
  boxsep=4pt,
  sharp corners,
]
\begin{center}
\textbf{\textcolor{blue}{Durée de l'exercice : 1h40 \\
Coefficient : 0.5 \\
Note sur 20 points. \\
Le travail est à réaliser en autonomie, mais de l'aide ponctuelle peut être apportée à d'autres élèves, à condition d'être validée par le professeur. Cette aide sera valorisée dans la note finale. Vous pouvez vous aider d'Internet pour réaliser votre devoir.
}}
\end{center}
\end{tcolorbox}

\vspace{0.3em}

\begin{center}
  \textbf{\textcolor{red}{Attention au temps - si vous êtes bloqués, avancez dans le devoir.}}
\end{center}


\section*{Exercice 1 - Somme des angles dans le triangle}

\begin{questions}

\question Sur Geogebra : 
\begin{parts}
  \part Tracer un triangle quelconque et afficher / mesurer ses trois angles.
  \part Faire la somme de ces mesures, en la stockant dans une variable appelée \textit{total\_angles\_1}.
  \part Tracer un autre triangle quelconque sur le même fichier, et refaire les mêmes étapes que précédemment, en stockant la somme des angles dans une variable appelée \textit{total\_angles\_2}. 
  \part Que constate-t-on ?
  \part Exporter le fichier Geogebra en le sauvegardant sur votre tablette, dans un dossier approprié, au format .ggb, et en le nommant \textit{Prénom Nom - Somme Angles Triangles.ggb}
\end{parts}
\end{questions}

\section*{Exercice 2 - A la recherche du cercle magique}

\begin{questions}

\question \textbf{Partie 1} : Sur papier, et en utilisant uniquement la règle et le compas, reproduire le triangle ci-dessous et essayer à l'aide du compas de construire un cercle qui passe par les trois sommets du triangle. Noter E son centre.


  \begin{figure}[H]
  \centering
  \includegraphics[width=0.5\textwidth]{activite_2_01.jpg}
  \caption{Exercice Geogebra}
\end{figure}

\question \textbf{Partie 2}
\begin{parts}
  \part Construire les médiatrices des trois côtés du triangle (on peut utiliser l'une ou l'autre des 2 méthodes connues)
  \part Que constate-t-on ?
\end{parts}

\question \textbf{Partie 3} : Les médiatrices se coupent en un point appelé point de concours.  
On dit que les trois médiatrices sont concourantes.
\begin{parts}
  \part Nommer O le point de concours.
  \part Tracer le cercle de centre O et passant par A.
  \part Que constate-t-on ?
\end{parts}

\question \textbf{Partie 4} : Le cercle tracé à la partie 3 s'appelle le cercle circonscrit au triangle ABC. Sur un nouveau fichier geogebra, 
\begin{parts}
  \part Tracer un triangle quelconque et tracer son cercle circonscrit.
  \part Tracer ensuite sur le même fichier un triangle rectangle et tracer son cercle circonscrit. Que constate-t-on ?
  \part Sauvegarder le fichier Geogebra en le nommant \textit{Prénom Nom - Cercle magique.ggb}
\end{parts}

\end{questions}


\section*{Rendu du devoir - Partie 1 - Envoi d'email}

\begin{questions}

\question Aller sur EcoleDirecte, et envoyer un email dont le titre est "Prénom Nom - Devoir Geogebra" et qui contient en pièce jointe le fichier GeoGebra de l'exercice 1 (\textit{Prénom Nom - Somme Angles Triangles.ggb}) et de l'exercice 2 (\textit{Prénom Nom - Cercle magique.ggb}). 
\question Puis coller sur une copie double les différents dessins et commentaires faits dans les 2 exercices. Vous inscrirez votre prénom et votre nom sur la copie double.
\end{questions}

\section*{Exercice 1 - Suite}

\textcolor{red}{Cette section fait partie intégrante du devoir. Je l'ai mise là parce que je veux être certain que vous avez pris le temps de m'envoyer l'email EcoleDirecte avant la fin de l'heure}

\begin{questions}
  
\question Sur une feuille de papier (quadrillée)
\begin{parts}
  \part Tracer un nouveau triangle "assez grand" et le découper avec précision en suivant son contour.
  \part Nommer ses trois angles $\hat{A}$, $\hat{B}$ et $\hat{C}$.

\begin{figure}[H]
  \centering
  \includegraphics[width=0.5\textwidth]{activite_1_01.jpg}
  \caption{Exercice Geogebra}
\end{figure}

  \part Réaliser un pliage afin de rabattre les trois sommets du triangle vers l'intérieur pour former un rectangle. Vous collerez votre figure sur votre copie double.

  \begin{figure}[H]
  \centering
  \includegraphics[width=0.5\textwidth]{activite_1_02.jpg}
  \caption{Exercice Geogebra}
\end{figure}

  \part Expliquer pourquoi ce pliage permet d'expliquer le résultat établi dans la partie 1.
\end{parts}

\end{questions}

\section*{Rendu du devoir - Partie 2 - Rendu de la copie}

\begin{questions}
\question Coller sur une copie double les différents dessins et commentaires faits dans les 2 exercices. Vous inscrirez votre prénom et votre nom sur la copie double. Rendre la copie double
\end{questions}

\vspace{2em}

\section*{Barème}
\renewcommand{\arraystretch}{1.7}

{
  %\setlength{\LTleft}{-2.3cm}  % déborde de 1.5 cm dans la marge de gauche
  \setlength{\LTright}{0pt}    % rien de plus à droite

\centering
\begin{longtable}{|L{0.3\textwidth}|L{0.7\textwidth}|c|}
\hline
\rowcolor{lavande}
\textbf{Critère évalué} & \textbf{Description} & \textbf{Points} \\[2pt]
\hline

% ---------------- Exercice 1 ----------------
  Exercice 1 — Angles dans le triangle (GeoGebra + feuille)
  & Construction correcte d'un triangle sur GeoGebra, mesure des trois angles et calcul de la somme.  
  Fichier GeoGebra enregistré et nommé correctement ("Prénom Nom - Triangle avec Angle").  
  Sur feuille : triangle tracé proprement, pliage réalisé, explication cohérente du lien avec la somme des angles.
  & 5 \\
\hline

% ---------------- Exercice 2 ----------------
  Exercice 2 — "À la recherche du cercle magique"
  & Reproduction correcte des triangles demandés.  
  Médiatrices tracées avec soin, point de concours identifié et nommé O.  
  Cercle de centre O passant par A correctement tracé.  
  Cercles circonscrits des triangles équilatéral et rectangle réalisés, avec constats rédigés clairement pour chaque cas.
  & 5 \\
\hline

% ---------------- Propreté ----------------
  Propreté et soin du travail
  & Écriture lisible et assez régulière.  
  Figures tracées aux instruments (règle, équerre, compas), pas de tracés à main levée.  
  Peu de ratures, feuilles propres, découpages et collages soignés.
  & 2 \\
\hline

% ---------------- Rendu ----------------
  Rendu du devoir (email + copie double)
  & Email envoyé via ÉcoleDirecte avec un objet clair du type "Prénom Nom - Devoir Geogebra".  
  Fichier GeoGebra joint au bon format (.ggb). Les noms des fichiers respectent bien le format demandé. 
  Copie double rendue avec prénom, nom, classe, et tous les dessins / réponses collés ou recopiés proprement.
  & 4 \\
\hline

% ---------------- Entraide et questions ----------------
  Entraide et qualité des questions posées
  & L'élève demande de l'aide de manière précise (par exemple : "Je n'arrive pas à construire les médiatrices" plutôt que "Je n'ai pas compris"). Ou arrive à faire efficacement des recherches sur Internet.
  Il ou elle peut aider un camarade en expliquant une méthode sans simplement donner la réponse.  
  Les questions posées montrent une vraie recherche de compréhension.
  & 4 \\
\hline

% ---------------- Total ----------------
\rowcolor{green!35}
 \textbf{Total} & \textbf{Note maximale possible} & \textbf{20} \\
\hline

\end{longtable}
}

\end{document}
